% Part: first-order-logic
% Chapter: syntax-and-semantics
% Section: terms-formulas

\documentclass[../../../include/open-logic-section]{subfiles}

\begin{document}

\olfileid{fol}{syn}{frm}

\olsection{项与\printtoken{P}{formula}}

一旦给定了某个一阶语言~$\Lang L$，我们就可以定义由~$\Lang L$ 的基本词汇构建的表达式。其中特别包括\emph{项}和\emph{公式}。

\begin{defn}[项]
\ollabel{defn:terms}
语言~$\Lang L$ 的\emph{项}集合~$\Trm[L]$ 归纳定义如下：
\begin{enumerate}
\item 每个变元都是一个项。
\item $\Lang L$ 的每个常项符号都是一个项。
\item 若 $f$ 是一个 $n$ 元函数符号且 $t_1$, \dots, $t_n$
  是项，则 $\Atom{f}{t_1, \ldots, t_n}$ 是一个项。
\tagitem{limitClause}{
  此外皆非项。}{}
\end{enumerate}
不含任何变元的项称为\emph{闭项}。
\end{defn}

\begin{explain}
常项符号在语言和项的定义中作为单独的一类符号出现，但它们本也可以被当作零元函数符号包含进来。那样的话，项定义中的第二条就可以省去了。我们只需将 $n = 0$ 时的 $\Atom{f}{t_1, \ldots, t_n}$ 理解为 $f$ 本身即可。
\end{explain}

\begin{defn}[公式]
\ollabel{defn:formulas}
语言~$\Lang L$ 的\emph{公式}集合~$\Frm[L]$ 归纳定义如下：
\begin{enumerate}
\tagitem{prvFalse}{$\lfalse$ 是一个原子公式。}{}

\tagitem{prvTrue}{$\ltrue$ 是一个原子公式。}{}

\item 若 $R$ 是~$\Lang L$ 的一个 $n$ 元谓词符号且 $t_1$, \dots,
  $t_n$ 是~$\Lang L$ 的项，则 $\Atom{R}{t_1,\ldots, t_n}$ 是一个
  原子公式。

\item 若 $t_1$ 与 $t_2$ 是~$\Lang L$ 的项，则 $\Atom{\eq}{t_1, t_2}$
  是一个原子公式。

\tagitem{prvNot}{若 $!A$ 是一个公式，则 $\lnot !A$ 是
  一个公式。}{}

\tagitem{prvAnd}{若 $!A$ 与 $!B$ 是公式，则 $(!A \land
  !B)$ 是一个公式。}{}

\tagitem{prvOr}{若 $!A$ 与 $!B$ 是公式，则 $(!A \lor !B)$
  是一个公式。}{}

\tagitem{prvIf}{若 $!A$ 与 $!B$ 是公式，则 $(!A \lif !B)$
  是一个公式。}{}

\tagitem{prvIff}{若 $!A$ 与 $!B$ 是公式，则 $(!A \liff !B)$
  是一个公式。}{}

\tagitem{prvAll}{若 $!A$ 是一个公式且 $x$ 是一个变元，
  则 $\lforall[x][!A]$ 是一个公式。}{}

\tagitem{prvEx}{若 $!A$ 是一个公式且 $x$ 是一个变元，
  则 $\lexists[x][!A]$ 是一个公式。}{}

\tagitem{limitClause}{此外皆非公式。}{}
\end{enumerate}
\end{defn}

\begin{explain}
项集和公式集的定义都是\emph{归纳定义}。本质上，公式集是通过无穷多个阶段构建的。在初始阶段，我们规定所有原子公式都是公式；这对应于定义中的前几条，即关于
\iftag{prvTrue}{$\ltrue$, }{}%
\iftag{prvFalse}{$\lfalse$, }{}%
$\Atom{R}{t_1,\dots,t_n}$ 和 $\Atom{\eq}{t_1,t_2}$ 的情形。
因此，“原子公式”指的就是任何具有这种形式的公式。

定义的其他条款给出了从已构造的公式出发构造新公式的规则。在第二阶段，我们可以用这些规则从原子公式构造公式。在第三阶段，我们从原子公式以及第二阶段得到的公式构造新的公式，依此类推。只有最终在某个阶段被构造出来的才是公式，此外皆非公式。
\end{explain}

按照惯例，我们将 $\eq$ 写在它的主目之间并省略括号：$\eq[t_1][t_2]$ 是
$\Atom{\eq}{t_1,t_2}$ 的缩写。此外，$\lnot \Atom{\eq}{t_1,t_2}$ 被缩写为
$\eq/[t_1][t_2]$。在书写由 $!B$, $!C$ 通过二元联结词~$\ast$ 构造的公式 $(!B \ast !C)$ 时，我们通常会省略最外面的一对括号，直接写作~$!B \ast !C$。

\begin{intro}
有些逻辑教材要求，要使
\iftag{prvEx}{$\lexists[x][!A]$ }{}%
\iftag{notprvEx,notprvAll}{}{和 }%
\iftag{prvAll}{$\lforall[x][!A]$ }{}%
算作
\iftag{notprvEx,notprvAll}{公式}{公式}，
变元~$x$ 必须在~$!A$ 中出现。
不作此要求也不会有任何问题，反而会使事情更简单。
\end{intro}

\begin{tagblock}{defNot,defOr,defAnd,defIf,defIff,defTrue,defFalse,defEx,defAll}
\begin{defn}
使用被定义的算子构造的公式应按以下方式理解：

\begin{tagenumerate}{defTrue,defFalse,defNot,defOr,defAnd,defIf,defIff,defEx,defAll}

\tagitem{defTrue}{$\ltrue$ 是
  \iftag{prvFalse}{$\lnot\lfalse$}{$(!A \lor \lnot !A)$（$!A$ 为某固定原子公式）} 的缩写。}{}

\tagitem{defFalse}{$\lfalse$ 是
  \iftag{prvTrue}{$\lnot\ltrue$}{$(!A \land \lnot !A)$（$!A$ 为某固定原子公式）} 的缩写。}{}

\tagitem{defNot}{$\lnot !A$ 是 $!A \lif \lfalse$ 的缩写。}{}

\tagitem{defOr}{$!A \lor !B$ 是
  \iftag{prvAnd}{$\lnot(\lnot !A \land \lnot !B)$}{$\lnot !A \lif
    !B$} 的缩写。}{}

\tagitem{defAnd}{$!A \land !B$ 是
  \iftag{prvOr}{$\lnot(\lnot !A \lor \lnot !B)$}{$\lnot (!A \lif
    \lnot !B)$} 的缩写。}{}

\tagitem{defIf}{$!A \lif !B$ 是
  \iftag{prvOr}{$\lnot !A \lor !B)$}{$\lnot (!A \land \lnot !B)$} 的缩写。}{}

\tagitem{defIff}{$!A \liff !B$ 是 $(!A \lif !B) \land (!B
  \lif !A)$ 的缩写。}{}

\tagitem{defAll}{$\lforall[x][!A]$ 是 $\lnot\lexists[x][\lnot !A]$ 的缩写。}{}

\tagitem{defEx}{$\lexists[x][!A]$ 是 $\lnot\lforall[x][\lnot !A]$ 的缩写。}{}
\end{tagenumerate}
\end{defn}
\end{tagblock}

当我们在针对特定应用的语言中工作时，通常会将二元谓词符号和函数符号写在相应的项之间，
例如，在算术语言中写 $t_1 < t_2$ 和 $(t_1 + t_2)$，在集合论语言中写 $t_1 \in t_2$。
算术语言中的后继函数甚至按惯例写在它的\emph{自变元之后}：$t'$。然而，从形式上说，
这些分别只是 $\Atom{\Obj A^2_0}{t_1, t_2}$、$\Obj f^2_0(t_1, t_2)$、$\Atom{\Obj A^2_0}{t_1, t_2}$ 以及 $f^1_0(t)$ 的约定缩写。

\begin{defn}[语形同一性]
符号 $\ident$ 表示符号串之间的语形同一性，即 $!A \ident !B$ 当且仅当 $!A$ 和 $!B$ 是具有相同长度的符号串，并且在每个位置上都包含相同的符号。
\end{defn}

符号 $\ident$ 的两侧可以是由并置运算得到的符号串，
例如 $!A \ident (!B \lor !C)$ 表示：符号串 $!A$ 与按顺序并置一个左括号、字符串 $!B$、符号 $\lor$、字符串 $!C$ 和一个右括号所得到的符号串是同一个。
在这种情况下，我们就知道 $!A$ 的第一个符号是左括号，$!A$ 包含 $!B$ 作为子串（从第二个符号开始），该子串后面跟着 $\lor$，等等。

由于项和公式是通过归纳定义从基本元素构建起来的，我们可以使用以下归纳原理来证明关于它们的性质。

\begin{lem}[\emph{项上的归纳原理}]
    \ollabel{lem:trmind}
    设 $\Lang L$ 为一阶语言。
    若某个性质~$P$ 满足：
    %
    \begin{enumerate}
        \item 它对每个变元~$v$ 成立，
        %
        \item 它对 $\Lang L$ 中的每个常项符号~$a$ 成立，并且
        %
        \item 只要它对 $t_1$,~\dots, $t_n$ 成立，它对 $f(t_1,\dotsc,t_n)$ 也成立（假设 $t_1$,~\dots, $t_n$ 是 $\Lang{L}$ 的项，且 $f$ 是 $\Lang L$ 的一个 $n$ 元函数符号），
    \end{enumerate}
    那么 $P$ 对 $\Trm[L]$ 中的每个项都成立。
\end{lem}

\begin{prob}
    证明 \olref[fol][syn][frm]{lem:trmind}。
\end{prob}

\begin{lem}[\emph{公式上的归纳原理}]
    \ollabel{thm:frmind}
    设 $\Lang L$ 为一阶语言。
    若某个性质~$P$ 对所有原子公式都成立，并且满足：
    %
    \begin{enumerate}
        \tagitem{prvNot}{只要它对 $!A$ 成立，则它对 $\lnot !A$ 也成立；}{}
        \tagitem{prvAnd}{只要它对 $!A$ 和 $!B$ 成立，则它对 $(!A \land !B)$ 也成立；}{}
        \tagitem{prvOr}{只要它对 $!A$ 和 $!B$ 成立，则它对 $(!A \lor !B)$ 也成立；}{}
        \tagitem{prvIf}{只要它对 $!A$ 和 $!B$ 成立，则它对 $(!A \lif !B)$ 也成立；}{}
        \tagitem{prvIff}{只要它对 $!A$ 和 $!B$ 成立，则它对 $(!A \liff !B)$ 也成立；}{}
        \tagitem{prvEx}{只要它对 $!A$ 成立，则它对 $\lexists[x][!A]$ 也成立；}{}
        \tagitem{prvAll}{只要它对 $!A$ 成立，则它对 $\lforall[x][!A]$ 也成立；}{}
  \end{enumerate}
  (假设 $!A$ 和 $!B$ 是 $\Lang{L}$ 的公式)，
  那么 $P$ 对 $\Frm[L]$ 中的所有公式都成立。
\end{lem}

\end{document}